
\section{这是译文}
\label{sec:introduction}




这是译文 \emph{这是译文} 这是译文 \emph{这是译文} 这是译文.

这是译文 \emph{这是译文} ($\sis$) 这是译文 \emph{这是译文} ($\lwe$) 这是译文~\cite{Ajtai96}, 这是译文 (这是译文) 这是译文 $\gapsvp$, 这是译文 (这是译文.,~\cite{DBLP:journals/siamcomp/MicciancioR07}), 这是译文 $\sis$ 这是译文 $\sis$ 这是译文~\cite{Ajtai96} 这是译文~\cite{DBLP:journals/eccc/ECCC-TR96-042}, 这是译文~\cite{DBLP:conf/crypto/MicciancioV03,DBLP:conf/pkc/Lyubashevsky08,DBLP:conf/asiacrypt/KawachiTX08}, 这是译文~\cite{DBLP:conf/stoc/GentryPV08,DBLP:conf/eurocrypt/CashHKP10,DBLP:conf/pkc/Boyen10,DBLP:conf/eurocrypt/MicciancioP12,DBLP:conf/eurocrypt/Lyubashevsky12}. 

这是译文: \[
\parens{\parens{ \veca_i, \inner{\vc{a}_i, \vc{s}}+e_i \bmod q}}_{i}
\quad \text{and} \quad
\parens{\parens{ \veca_i, u_i }}_{i}~,
\] 这是译文 $\vecs$ 这是译文 $\Z_q^n$ 这是译文 $\veca_i \in \Z_q^n$,  $u_i$ 这是译文 $\Z_q$, 这是译文'' $e_i \in \Z$ 这是译文 (这是译文) 这是译文 $\alpha q$ 这是译文 $\alpha = o(1)$.
%\dnote{I replaced $\ll$ by $o(\cdot)$}


这是译文 $\lwe$ 这是译文~\cite{DBLP:journals/jacm/Regev09,DBLP:conf/crypto/PeikertVW08,DBLP:conf/ctrsa/LindnerP11} 这是译文~\cite{DBLP:conf/stoc/PeikertW08,DBLP:conf/stoc/Peikert09,DBLP:conf/eurocrypt/MicciancioP12} 这是译文~\cite{DBLP:conf/crypto/PeikertVW08}, 这是译文~\cite{DBLP:conf/stoc/GentryPV08,DBLP:conf/eurocrypt/CashHKP10,DBLP:conf/eurocrypt/AgrawalBB10,DBLP:conf/crypto/AgrawalBB10}, 这是译文 (这是译文., \cite{DBLP:conf/tcc/AkaviaGV09,DBLP:conf/crypto/ApplebaumCPS09,DBLP:conf/innovations/GoldwasserKPV10}), 这是译文~\cite{DBLP:conf/focs/BrakerskiV11,DBLP:conf/innovations/BrakerskiGV12,DBLP:conf/crypto/Brakerski12} (这是译文~\cite{DBLP:conf/stoc/Gentry09}), 这是译文~\cite{DBLP:conf/focs/KlivansS06}.

这是译文 $\sis$ 这是译文 $\lwe$ 这是译文 $\lwe$~\cite{DBLP:journals/jacm/Regev09} 这是译文 $\sis$ 这是译文 \emph{这是译文}. 这是译文 $\lwe$, 这是译文 (这是译文) 这是译文 \emph{这是译文} 这是译文 \emph{这是译文} 这是译文 $\lwe$.

这是译文~\cite{DBLP:conf/stoc/Peikert09} (这是译文~\cite{DBLP:conf/crypto/LyubashevskyM09}).  这是译文 $\lwe$ 这是译文 \emph{这是译文} 这是译文 $\lwe$ 这是译文~\cite{DBLP:conf/focs/KlivansS06}, 这是译文 (这是译文) 这是译文 $\lwe$ 这是译文 $2$ (这是译文), 这是译文~\cite{DBLP:conf/stoc/Peikert09} 这是译文 $\lwe$ 这是译文 $\lwe$ 这是译文.


这是译文 $\lwe$ 这是译文 \emph{这是译文} 这是译文 $\lwe$ 这是译文.

\myparagraph{Main result} 这是译文 $\lwe$ 这是译文 $\lwe$ 这是译文.

\onote{should we add somewhere a formal statement of our main result? at least a combination of sections 3 and 4?}
\begin{theorem}[这是译文]\label{thm:informal} 这是译文 $n$-这是译文 $\lwe$ 这是译文 $\poly(n)$ 这是译文 $\sqrt{n}$.
\end{theorem}
这是译文 $\lwe$ 这是译文.

\myparagraph{Techniques} 这是译文.

这是译文 $\lwe$ 这是译文 $q$ 这是译文 $n$ 这是译文 $\lwe$ 这是译文 (这是译文) 这是译文 $p=\poly(n)$ 这是译文 $n \log_2 q$. 这是译文~\ref{thm:informal} 这是译文~\cite{DBLP:conf/stoc/Peikert09}, 这是译文 $q=2^n$ 这是译文 $n$-这是译文 $\gapsvp$. 这是译文 $n$ 这是译文 $n \log_2 q$ 这是译文 (这是译文).%invariant

这是译文 $\Z_q$ 这是译文 $\Z_p$ 这是译文 $a \in \{0,\ldots,q-1\}$ 这是译文 $\floor{pa/q} \in \{0,\ldots,p-1\}$. 这是译文 $\lwe$, 这是译文 (这是译文~\ref{sec:modexpansion}). 这是译文 $\lwe$ 这是译文 $\vecs$, 这是译文 $\vecs$ 这是译文 (这是译文~\ref{sec:shortsecret}) 这是译文 $\lwe$ (这是译文 $n$) 这是译文 $\Z_q^n$ 这是译文 $\lwe$ (这是译文 $n \log_2 q$) 这是译文 $\{0,1\}$. 这是译文 (这是译文~\ref{sec:theshortsecretreduction}) 这是译文-$\lwe$'' 这是译文 $\lwe$ 这是译文 (这是译文~\ref{sec:extlwe}).


这是译文~\cite{DBLP:conf/focs/BrakerskiV11,DBLP:conf/innovations/BrakerskiGV12,DBLP:conf/crypto/Brakerski12}. 这是译文 $\lwe$ 这是译文.


%
这是译文 $\lwe$ 这是译文 $\{0,1\}$ 这是译文.~\cite{DBLP:conf/innovations/GoldwasserKPV10} 这是译文 $\lwe$ 这是译文 $\{0,1\}$-$\lwe$ 这是译文~\cite{DBLP:conf/innovations/GoldwasserKPV10} (这是译文), 这是译文 $\{0,1\}$-$\lwe$ 这是译文, $\lwe$ 这是译文 $\{0,1\}$-$\lwe$ 这是译文.


这是译文~\cite{DBLP:conf/innovations/GoldwasserKPV10} 这是译文 $\lwe$ 这是译文 $\vecz$ 这是译文~\cite{DBLP:conf/crypto/ONeillPW11} 这是译文 (这是译文)这是译文~\cite{DBLP:conf/pkc/Alperin-SheriffP12} 这是译文.



\myparagraph{Broader perspective} 这是译文~\ref{thm:informal}, 这是译文 $\lwe$. 这是译文~\ref{sec:modexpansion} 这是译文 $\lwe$ 这是译文 $n$ 这是译文 $p$ 这是译文 $\lwe$ 这是译文 $n/k$ 这是译文 $p^k$ (这是译文~\ref{cor:mod-dim-tradeoff}). 这是译文 $n$-这是译文 $\lwe$ 这是译文 $q$ 这是译文 $n \log_2 q$. 这是译文 $n$ 这是译文 $q$ 这是译文 $n \log_2 q$ 这是译文 $\lwe$.

这是译文 (这是译文 $n \log_2 q$ 这是译文), 这是译文 $n$-这是译文 $\lwe$ 这是译文 $2^n$ 这是译文 $n^2$-这是译文 $\lwe$ 这是译文, $n$-这是译文 $\lwe$ 这是译文 $2^n$, 这是译文~\cite{DBLP:conf/stoc/Peikert09} 这是译文 $n$-这是译文 $n^2$-这是译文 $\exp(\widetilde\Omega(n^2))$ 这是译文~\cite{DBLP:conf/stoc/Peikert09} (这是译文~\ref{thm:informal}) 这是译文~\cite{DBLP:journals/jacm/Regev09}.

\onote{for the final version, consider adding precise statements for this and the next point} 这是译文 $1$-这是译文 $\lwe$ 这是译文 $2^n$ 这是译文 $n$-这是译文 $\lwe$ 这是译文 $1$-这是译文 $\lwe$ 这是译文~\cite{DBLP:conf/crypto/BonehV96}. 这是译文~\cite{DBLP:conf/stoc/AjtaiD97} 这是译文~\cite{DBLP:journals/jacm/Regev04}, 这是译文~\cite{RegevSurvey}. 这是译文 $1$-这是译文-$\lwe$ 这是译文~\cite{DBLP:conf/eurocrypt/LyubashevskyPR10} (这是译文~1, 这是译文~$\Z$). 这是译文~\cite{DBLP:conf/scn/GentryHPS12}, 这是译文-$\lwe$, 这是译文 (这是译文). 这是译文-$\sis$ 这是译文~\cite{DBLP:conf/icalp/LyubashevskyM06,DBLP:conf/tcc/PeikertR06} 这是译文 $\lwe$-这是译文-$\sis$ 这是译文. (这是译文 $1$, 这是译文.)

这是译文 (这是译文) \lwe 这是译文~\ref{cor:mod-reduction-2}.  这是译文 (这是译文.,~\cite{DBLP:journals/jacm/Regev09,DBLP:conf/stoc/Peikert09,DBLP:conf/crypto/MicciancioM11,DBLP:conf/eurocrypt/MicciancioP12}) 这是译文 \emph{这是译文}, 这是译文.

\onote{for final version: decide what to do with real-LWE}

\onote{for final version, consider applying these ideas to SIS, especially the 1-dim case which is like subset sum, and more generally, modulus/dimension tradeoff for SIS. Chris said:
The idea is to do a similar transformation as in Lemma 4.4, but on SIS samples, going from a in $\Z_q^n$ to $a' \in \Z_{q'}^{n'}$.  Then given a good SIS solution $x$ to $A'x = 0$, we also have $G^T A' x = 0$ and therefore $A x \approx 0$.  Thus $x$ is a solution to the original SIS instance, because it suffices to find a preimage of "nearly" 0.  This is because if $Ax = e \approx 0$, then we get the homogeneous solution $[I | A] (e; x) = 0$, and the instance $[I | A]$ is just SIS in (hermite) normal form.}

\dnote{Something that I find interesting is that the complexity of the
  best known attack is~$2^{O(n \log q/ \log^2 \alpha)}$, which is
  constant wrt $n \log q$. This was not reflected by our previous
  understanding of the complexity of \lwe, and we manage to close the
  gap between the hardness results and the algorithms. By the way,
  this~$2^{O(n \log q/ \log^2 \alpha)}$ bound suggests it may be
  possible to transfer some of~$n$ and~$\log q$ into~$\log \alpha$. Do
  we know \lwe self-reductions from, say, $q,\alpha$ to~$p,\beta$ such
  that~$\log q/ |\log \alpha| = \log p/|\log \beta|$? or from
  $n,\alpha$ to~$n',\beta$ such that~$n/ |\log \alpha| = n'/|\log
  \beta|$}

\onote{maybe mention Ajtai's 2005 paper using Dirichlet's problem?}

\myparagraph{Open questions} 这是译文~\ref{thm:informal} 这是译文~\cite{DBLP:conf/stoc/Peikert09} 这是译文~\cite{DBLP:journals/jacm/Regev09} 这是译文 $\sis$.  这是译文~\cite{DBLP:conf/stoc/Peikert09} 这是译文~\cite{DBLP:journals/jacm/Regev09} 这是译文 $\sis$ 这是译文~\ref{thm:informal}  这是译文~\cite{DBLP:conf/stoc/Peikert09} 这是译文 \gapsvp (这是译文 \bdd 这是译文 \usvp~\cite{DBLP:conf/crypto/LyubashevskyM09}), 这是译文 \sivp, 这是译文~\cite{DBLP:journals/jacm/Regev09} 这是译文 \sis 这是译文-\lwe, 这是译文 \gapsvp 这是译文. 

这是译文~\cite{DBLP:conf/stoc/Peikert09} 这是译文~\ref{thm:informal}. 这是译文 $n$-这是译文 $\lwe$ 这是译文 $n$-这是译文 (这是译文 $\sqrt{n}$-这是译文)? 这是译文~\cite{DBLP:journals/jacm/Regev09}.

这是译文 (这是译文) 这是译文 $\lwe$ 这是译文 $2^{O(\sqrt{n})}$. 这是译文~\ref{thm:informal}, 这是译文 $2^{O(\sqrt{n})}$-这是译文 \emph{这是译文} 这是译文~\cite{DBLP:journals/jacm/Regev09}. 这是译文 $2^{O(\sqrt{n})}$-这是译文~\cite{DBLP:journals/siamcomp/Kuperberg05} (这是译文~\cite{DBLP:journals/siamcomp/Regev04}).




%%% Local Variables:
%%% mode: latex
%%% TeX-master: "lwehard"
%%% End:
